\documentclass[11pt, letterpaper]{article}

% --- Packages ---
\usepackage[margin=0.85in]{geometry}
\usepackage{hyperref}
\usepackage{titlesec}
\usepackage{xcolor}
\usepackage[T1]{fontenc}
\usepackage{mathpazo} % Palatino font for classic academic elegance
\usepackage{microtype} % Superior typography and letter spacing
\usepackage{enumitem} % Advanced list control

% --- Colors ---
\definecolor{primary}{RGB}{0, 0, 0}
\definecolor{linkcolor}{RGB}{25, 84, 166}

% --- Hyperlinks ---
\hypersetup{
    colorlinks=true,
    linkcolor=linkcolor,
    filecolor=linkcolor,
    urlcolor=linkcolor,
    pdftitle={Academic CV - Xiyuan GUO},
}

% --- Section Formatting ---
\titleformat{\section}{
  \vspace{6pt}\Large\scshape\raggedright\bfseries\color{primary}
}{}{0em}{}[\color{primary}\titlerule \vspace{6pt}]

% --- Custom CV Environments ---
\newenvironment{cvitemize}{
    \begin{itemize}[leftmargin=0.2in, label={\textbullet}, itemsep=4pt, parsep=0pt, topsep=4pt]
}{
    \end{itemize}
}

\newcommand{\cvheading}[4]{
    \vspace{4pt}
    \noindent\textbf{#1} \hfill #2 \\
    #3 \hfill \textit{#4} \vspace{2pt}
}

\begin{document}

% --- Heading ---
\begin{center}
    {\Huge \scshape \textbf{Xiyuan (Tommy) GUO}} \\ \vspace{8pt}
    B.A. Candidate in Computer Science, University of Virginia \\ \vspace{4pt}
    \href{mailto:mhn4xq@virginia.edu}{mhn4xq@virginia.edu} $\vert$ (516) 474-3279 $\vert$ \href{https://xguo.info}{xguo.info} $\vert$ \href{https://github.com/xguot}{github.com/xguot}
\end{center}
\vspace{8pt}

% --- Scientific Focus ---
\section{RESEARCH FOCUS}
\noindent Convex optimization; computational imaging; missing-data imputation; quantum error correction.

% --- Education ---
\section{EDUCATION}
\cvheading{University of Virginia}{Expected May 2027}{B.A. in Computer Science}{Charlottesville, VA}
\begin{cvitemize}
    \item \textbf{GPA:} 3.73/4.00
    \item \textbf{Relevant Coursework:} Graduate Quantum Computing (PHYS 5880, in progress).
    \item \textbf{Independent Study:} Quantum Information (Preskill lecture series); Quantum Mechanics (Griffiths, Sakurai).
\end{cvitemize}
\cvheading{Kimball Union Academy}{2024}{High School Diploma}{Meriden, NH}

% --- Grants & Fellowships ---
\section{GRANTS AND FELLOWSHIPS}
\begin{cvitemize}
    \item \textbf{Double Hoo Research Grant}, University of Virginia \hfill \textit{2026} \\
    \small{Secured competitive grant funding to drive independent, interdisciplinary computational research alongside graduate mentorship.}
\end{cvitemize}

% --- Publications & Manuscripts ---
\section{PUBLICATIONS \& MANUSCRIPTS IN PREPARATION}
\begin{cvitemize}
    \item Guo, X., [Co-authors], Ke, L. "Transition-metal magnetic anisotropy and its relationship to structural motifs." \textit{[Manuscript in preparation]}.
    \item Guo, X., [Co-authors], Liu, M. "Joint Spectro-Polarimetric Lensless Reconstruction via Conic-Constrained ADMM." \textit{[Manuscript in preparation]}. Code available on request.
    \item Guo, X., [Co-authors], Tong, X.C. "PRISM: Lagrangian Projection Engines for Multivariate Longitudinal Missing Data Imputation." \textit{[Manuscript in preparation]}. Code: \href{https://github.com/xguot/prism}{github.com/xguot/prism}.
\end{cvitemize}

% --- Research Experience ---
\section{PROFESSIONAL RESEARCH EXPERIENCE}

\cvheading{Research Assistant}{Feb 2026 -- Present}{Yale University (Prof. Mengxia Liu's Group)}{New Haven, CT}
\begin{cvitemize}
    \item Built a joint spectro-polarimetric forward model recovering a 324-channel datacube (spatial $\times$ 81 wavelengths $\times$ 4 Stokes parameters) from a single multiplexed lensless speckle image.
    \item Formulated the reconstruction as a conic-constrained inverse problem: enforcing a Lorentz cone constraint on the Stokes vector ($S_0 \geq \|S_{1:3}\|$) to guarantee physical polarization states.
    \item Implemented a consensus ADMM solver in PyTorch with $L_1$ spectral sparsity and Sherman--Morrison linear solves; runs on GPU via SLURM.
    \item Calibrated PSF libraries using Gram--Schmidt orthogonalization to remove channel cross-talk; spectral unmixing via non-negative least squares.
\end{cvitemize}

\cvheading{Research Assistant}{Oct 2025 -- Present}{UVA (Prof. Liqin Ke's Group)}{Charlottesville, VA}
\begin{cvitemize}
    \item Used Density Functional Theory (DFT) to study transition-metal magnetic anisotropy and its relationship to structural motifs in crystal lattices.
    \item Built interactive 3D atomic-structure visualizations (Three.js, JSmol) for real-time structural inspection.
    \item Developed a shared ML pipeline (Go, Next.js) for collaborators to run predictive models on magnetic-property datasets.
    \item Designing RareEarth DB, a rare-earth analogue of the Materials Project: 4,648 curated anisotropy measurements linked to ICSD/Pearson structures, with Fireworks-orchestrated DFT workflows on UVA HPC; DOE funding proposal in preparation.
\end{cvitemize}

\cvheading{Research Assistant}{Apr 2026 -- Present}{UVA BAL Lab (Prof. Xin Cynthia Tong's Group)}{Charlottesville, VA}
\begin{cvitemize}
    \item Developed PRISM, a constrained optimization framework for multivariate longitudinal missing-data imputation: projects initial imputations onto the model-implied covariance structure to correct the variance attenuation of standard conditional-expectation methods.
    \item Built the solver in C++17 (RcppArmadillo): projected gradient descent with Armijo line search, exact mean-constraint projection, and KKT convergence checks.
\end{cvitemize}

\cvheading{Research Intern}{May 2026 -- Aug 2026}{UVA Biocomplexity Institute}{Charlottesville, VA}
\begin{cvitemize}
    \item Engineered a Julia simulation engine for Independent Cascade diffusion and trained PyTorch Geometric GCNs to estimate total epidemic cascade size from sparse, partially observed network infections.
    \item Developed an agentic workflow engine (Pydantic AI, NetworkX) translating natural-language queries into verified execution DAGs for network-analysis pipelines, with deterministic validation and graph deduplication cutting redundant computation by $>40\%$.
\end{cvitemize}

% --- Open-Source Software ---
\section{OPEN SOURCE SOFTWARE}
\begin{cvitemize}
    \item \textbf{stim:} Authored PR \#1099 to the quantum stabilizer circuit simulator (Google; accepted in review).
    \item \textbf{PRISM:} C++17 constrained-optimization engine for longitudinal missing-data imputation (see research above).
    \item \textbf{qudec:} Simulation toolkit for quantum LDPC decoding (BP+OSD, LP, ADMM); characterized BP failure modes on degenerate codes and benchmarked decoders under phenomenological noise.
    \item \textbf{sanjaya:} Academic literature extraction pipeline (OpenAlex API, Scrapy, Playwright) for structured bilingual dataset export.
    \item \textbf{difi:} Go-based Git diff TUI; Homebrew Core, 365+ GitHub stars, \#5 on Hacker News.
\end{cvitemize}

% --- Awards ---
\section{HONORS AND AWARDS}
\begin{cvitemize}
    \item \textbf{Double Hoo Research Grant}, University of Virginia, 2026
    \item \textbf{Dean's List of Distinguished Students}, University of Virginia
    \item \textbf{Silver Division}, USA Computing Olympiad (USACO)
\end{cvitemize}

% --- Skills ---
\section{TECHNICAL SKILLS}
\begin{cvitemize}
    \item \textbf{Math \& Physics:} Advanced Theoretical Quantum Mechanics, Quantum Information Theory, Convex Optimization, Matrix Calculus, Linear Algebra.
    \item \textbf{Machine Learning \& Computing:} PyTorch/CUDA, JAX, Scikit-Learn, Pandas, RcppArmadillo.
    \item \textbf{Languages:} Python, C/C++, Julia, Rust, Go, TypeScript, R, SQL, Bash, \LaTeX.
    \item \textbf{Infrastructure \& Cloud:} Linux (openSUSE/RHEL), Docker, K3s Kubernetes, Slurm HPC, Git.
\end{cvitemize}

% --- Service ---
\section{COMMUNITY SERVICE}
\cvheading{Volunteer Firefighter}{2022 -- 2024}{Plainfield Fire Department}{Plainfield}
\begin{cvitemize}
    \item Responded to over 200 emergency medical, fire, and rescue calls annually, operating effectively in high-stress environments to provide critical incident mitigation and community support.
\end{cvitemize}

\end{document}
